\section{Method}

H001 is a \method{} framework for 3D scene graph relation predictions. It is
not a replacement relation predictor and not a standalone verifier script. The
framework is built around four design requirements: preserve relation-row
identity, join explicit object-pair geometry, calibrate geometry validity, and
evaluate recall and violation jointly. Figure~\ref{fig:framework} presents
these requirements as a pipeline from relation-source outputs to calibrated
operating points.

\subsection{Relation-Row Standardization}

Different relation sources expose predictions in different formats. We convert
each source output into a shared row contract:
\begin{equation}
(\mathrm{scan},\mathrm{subgraph},s,o,p,a,\mathrm{score}^{sem},\mathrm{source}).
\end{equation}
This row is the unit of comparison across VL-SAT, Open3DSG, controls, and
future optional sources. The standardized row keeps the semantic score from the
source model but adds the fields needed for identity-preserving geometry joins
and denominator-stable evaluation.

\subsection{Identity-Preserving Geometry Join}

Geometry is joined by $(\mathrm{scan},\mathrm{subgraph},s,o)$, not by category
names or unordered object-pair statistics. This matters because the same
predicate can be plausible for many object categories while being false for a
particular instance pair. The wrong-pair and shuffled-geometry controls
directly test this design choice by breaking or perturbing object-pair identity
while preserving parts of the geometry distribution.

\subsection{Family-Specific Geometry Checks}

H001 uses relation-family-specific checks for \texttt{support\_contact},
\texttt{proximity}, and \texttt{relative\_vertical}. For support/contact, the
evidence includes point/local contact, vertical support, OBB overlap, and
support subtypes. For proximity, the evidence includes object distance and pair
scale. For relative vertical relations, the evidence includes centroid and
extent ordering along the vertical axis. The verifier returns violated only
when the evidence is strong enough; otherwise ambiguous or weak geometry can be
marked uncertain.

\subsection{Calibrated Geometry Validity and Re-Ranking}

The probabilistic operating point uses a frozen calibrator:
\begin{equation}
\pgeom_i = C_a(\phi(g_i)),
\end{equation}
where $\phi(g_i)$ is the geometry feature vector and $C_a$ is either the pooled
or family-specific calibrator trained on train-dev calibration positives and
counterfactual negatives. The main probabilistic re-ranking score is
\begin{equation}
\mathrm{score}^{prob}_i = \mathrm{score}^{sem}_i \cdot \pgeom_i.
\end{equation}

We report four main operating points. \texttt{semantic\_only} ranks by the
source semantic score and serves as the reproduced source baseline.
\texttt{probabilistic\_recalibrated} ranks by the product of semantic
confidence and calibrated geometry validity.
\texttt{rule\_verified\_point\_subtype} removes hard violations before ranking
and is reported as a zero-violation diagnostic.
\texttt{family\_specific\_p\_geom\_valid} uses family-specific calibration and
provides a stricter violation-first operating point.

\subsection{Controls}

The control conditions test whether the result can be explained by simpler
signals. Geometry-only ranking tests whether semantic confidence remains
necessary. Distance-only ranking tests whether a generic spatial heuristic
explains the result. Shuffled-geometry and wrong-pair-geometry controls test
whether relation-level object-pair identity is necessary. If these controls
perform worse than the calibrated setting, the result supports the framework
design rather than a generic ``add geometry'' explanation.

